\subsection{Interview b1-aug05-1650}\label{int:b1-aug05-1650}
\subsubsection*{Metadata}
\begin{itemize}
  \item \textbf{ID}: b1-aug05-1650
  \item \textbf{Date/time}: 2026-08-05T16:50
  \item \textbf{Method}: in person
  \item \textbf{Language}: English
  \item \textbf{Location}: Slovakia
  \item \textbf{Interviewee}: age 44, other, Slovak (see notes — no notes section present in source) nationality, main language Slovak
  \item \textbf{Questions / Answers}: 3 / 3
  \item \textbf{Total word count}: 273
  \item \textbf{Corrections log}: \texttt{cleaned/b1-aug05-1650/corrections.log}
  \item \textbf{Notes}: Nationality field says '(see notes)' in the source but no notes section exists there — left verbatim, flagged as a data-quality gap.
\end{itemize}
\subsubsection*{Transcript}
INTERVIEW — Aug 5, 16:50

Interviewee: 44 years old, gender "other", main language Slovak, lives in Slovakia,
nationality Slovak (see notes).

---

Interviewer: So, my first question is: when you hear the word cult, what comes to your mind?

Interviewee: Tomato cult! [laughs] Okay — it's a local community of people growing local tomatoes.

Interviewer: If you think about them, what makes it a cult for you?

Interviewee: Exactly — dedication and commitment to their voluntary hobby, and the love and passion for what they're doing — which is getting back to roots. In these times, people in a big city are growing tomatoes, and they're changing varieties, and they're measuring the sugar [content] of the tomatoes, and they're kind of scientifically looking at growing the tomatoes, and the types, and — yeah.

Interviewer: Um, so, which features, then, for you, are really the most important to call this group a cult? It's a group?

Interviewee: It's a group, but without any kind of hierarchy — key members, or a set number of members. And it started as a group in the pub, and it got bigger on a Facebook group page. Now, after several years, we added an event that unites these people. Yeah, and the one word that would kind of describe it, in the sense of "cult" — is maybe kind of sarcastically named "the cult," because, you know, there's no reason why this should be a cult. But, yeah, maybe, like — I would say sarcastically, in the thinking of "fuck cult," or something like this. Okay, so, that's my opinion — I mean, I'm not sure, yeah.


---


\subsection{Interview b1-aug06-1955}\label{int:b1-aug06-1955}
\subsubsection*{Metadata}
\begin{itemize}
  \item \textbf{ID}: b1-aug06-1955
  \item \textbf{Date/time}: 2026-08-06T19:55
  \item \textbf{Method}: in person
  \item \textbf{Language}: English
  \item \textbf{Location}: Germany
  \item \textbf{Interviewee}: age 26, female, Ukrainian nationality, main language Ukrainian
  \item \textbf{Questions / Answers}: 4 / 7
  \item \textbf{Total word count}: 410
  \item \textbf{Corrections log}: \texttt{cleaned/b1-aug06-1955/corrections.log}
\end{itemize}
\subsubsection*{Transcript}
INTERVIEW — Aug 6, 19:55

Interviewee: 26 years old, female, main language Ukrainian, speaking in English, lives in Germany, Ukrainian nationality.


---

Interviewer: Um, when you hear the word cult, what comes to your mind?

Interviewee: Actually — the first thing I thought was this AI thing you told me about, this AI cult. But because I know your background and probably your thoughts.

Interviewer: Well, what's the first thing you think about, about cults?

Interviewee: I don't know, I don't really know a lot about those. [You mean] the first thing that came to your mind is this AI cult I talked about — do you remember?

Interviewer: Yes, I just remembered

Interviewee: You explained the theory that there exists some cult that says artificial intelligence could be like a god. Yes, you know — so, this is my first thought. Okay, yeah, but I don't actually really know about cults. Maybe I saw something in American Horror Story — there was some series, okay, I think it was, like, one season or something — and it was a cult where, uh, people were killing other people, dressed like clowns.

Interviewee: Yes — this is the first thing I think about, because I'm not — [laughs] sorry, sorry. Yeah, yeah. And so, in this American Horror Story season — this cult — what did it look like, and what made it a cult, do you think?

Interviewee: Uh, there was one guy who was like an icon for them, and he had, like, really good energy, I would say. He influenced them — his character, probably, his energy — and he was the main person who, like, helped other people to open up their bad side, you know. And, after that, they — I think — they were wearing clown costumes, and they were killing people. Yeah, so this is what I remember about cults, okay.

Interviewee: And Charles Manson! This is the guy also, and it's a similar thing — he has, like, a huge, great character and energy. He never actually killed real people with his own hands, but he influenced other people to do that.

Interviewer: And I kind of — yeah, it's interesting. So for you, the main component is this influence thing.

Interviewee: Yeah, that's like a strong concept, you know. I'm sometimes really interested in that — how it's possible to influence people this way, to create a huge cult like that.

---


\subsection{Interview b1-aug08-1420}\label{int:b1-aug08-1420}
\subsubsection*{Metadata}
\begin{itemize}
  \item \textbf{ID}: b1-aug08-1420
  \item \textbf{Date/time}: 2026-08-08T14:20
  \item \textbf{Method}: in person
  \item \textbf{Language}: English
  \item \textbf{Location}: Germany
  \item \textbf{Interviewee}: age 33, female, dual Australian/Iranian nationality, main language English
  \item \textbf{Questions / Answers}: 6 / 6
  \item \textbf{Total word count}: 466
  \item \textbf{Corrections log}: \texttt{cleaned/b1-aug08-1420/corrections.log}
\end{itemize}
\subsubsection*{Transcript}
INTERVIEW — Aug 8, 14:20

Interviewee: 33 years old, female, main language English, speaking in English, lives in
Germany, dual Australian/Iranian nationality.


---

Interviewer: Okay, perfect. So, first question — when you hear the word cult, what comes to your mind?

Interviewee: Oh, okay, yeah — not good. You need more words?

Interviewer: Um, yeah — well, like, maybe think of an example, or an image, or, like, what it means, kind of thing — whatever you picture.

Interviewee: Huh, okay. So, I'd picture something American, okay — brainwashing, suppressive, abusive. Um, yeah, usually white dudes. Yeah — like, white American dudes.

Interviewer: Any particular example?

Interviewee: Oh, I — yeah, I don't have any names, or — yeah, I don't remember any names.

Interviewer: That's okay. So, these attributes that you describe — for example, the white dudes — like, why would you associate it with white people?

Interviewee: I guess because I've seen documentaries, or media, and it was usually in those instances that I did see it — it was those kinds of people.

Interviewer: Okay, so — and then, yeah, interesting — so that's all the things that you've seen just from media. And so, these things that you saw in documentaries — like, why do you think they characterize them as cults?

Interviewee: I guess there's a definition of cult, and hopefully they're assessed based on some definition, but I guess a cult is anything that becomes too extreme in an ideology, and kind of just messes with the people themselves. They don't necessarily outwardly, like, do things — they just manipulate and control the people within the cult, and they just make their own little society. And, yeah, that's my understanding of cult.

Interviewer: Okay. And, like, given these things — there's still no examples that come to your mind when you think about it?

Interviewee: Yeah, I don't have any names. And, like, yeah, I kind of just think, "okay, I don't want that information in my brain anymore."  Yeah, so I guess I also see that often — in examples I've seen but don't remember the names of — religion, or some sort of spirituality, was used and then, like, exaggerated. Yeah, so I was curious, thinking about it myself, like — if someone is considered to be a Zionist, is it because they have the core belief that they want to annihilate the people of Palestine, and everyone else that they deem is putting them at risk? Then, why couldn't that also be considered a cult — because they are also raised and brainwashed from a young age, given this narrative, so their reality is completely different — which is kind of similar to what happens in cults, where children are raised with a completely different reality.


\subsection{Interview b1-aug08-2128}\label{int:b1-aug08-2128}
\subsubsection*{Metadata}
\begin{itemize}
  \item \textbf{ID}: b1-aug08-2128
  \item \textbf{Date/time}: 2026-08-08T21:28
  \item \textbf{Method}: in person
  \item \textbf{Language}: English
  \item \textbf{Location}: United States
  \item \textbf{Interviewee}: age 29, female, German nationality, main language German
  \item \textbf{Questions / Answers}: 6 / 6
  \item \textbf{Total word count}: 564
  \item \textbf{Corrections log}: \texttt{cleaned/b1-aug08-2128/corrections.log}
\end{itemize}
\subsubsection*{Transcript}
INTERVIEW — Aug 8, 21:28

Interviewee: 29 years old, gender Female, main language
German, speaking in English, lives in the US, nationality German.

---

Interviewer: Okay, great. So, my first question is: when you hear the word cult, what comes to your mind?

Interviewee: Something suspicious, something dangerous, something not normal. And also something a little bit frightening.

Interviewer: Okay, do you have a specific example that would come to your mind immediately?

Interviewee: I don't have a specific name, but it's more like — those yoga cults which existed some time ago. I don't have other names.

Interviewer: Okay, what is it that's "culty" to you about them?

Interviewee: I would say it's people participating in something, and then there's a person who has another agenda and kind of manipulates them, and uses, like — their, I don't know — participation, or — I think that's not the right word — but, like, their willingness, and, like, their good [nature] to participate in something. And just, like — yeah, I don't know — I need another word in English — but, like, what's the word in German? Um — "Missbrauch" [German for "abuse"/"misuse"]. Yeah, exactly — he's famous for abusing the trust — the trust, the consent, of the people — and using it for his own purposes.

Interviewer: Okay, if you were to, then, try to give a definition of the cult, which aspect would you feel is the most important of all the things you already said? Or maybe there are other ones that come to your mind?

Interviewee: I think it's benefiting from something the other person, or other people, aren't actually aware that is your intention. I think that, for me, is the core of it. Yeah, because when you have a religion, I don't know — you're all having the same thoughts, beliefs, for one thing, and it's like the common knowledge more, it is, okay — whereas with a cult, it's more — um, like, my opinion, it's — I don't know — someone is [misusing, "Missbrauch"] — [?heavily garbled stammering — see notes] — I think that it's the misuse, the distortion, the distortion, the disrespect — like, you're not doing what you should, what you're intending to do, and you're not doing what other people are expecting you to do — from, like, a moral, spiritual point of view, or — yes, ethically, ethically, morally, I think, all these. They're not really giving you something, they just pretend to give you something, and then you, being part of the cult  are not aware of that. You're being part of, like, a cult — you just think you're part of a community, or something like that. Yeah, so you don't actually know you're part of a cult — but, zooming out, you see both kind of perspectives, you see that. And there are some yoga communities doing this.

Interviewer: Yes, interesting. Which documentary do you remember?

Interviewee: I don't know, there were some Netflix documentaries — okay, there was one famous one in the U.S., but I don't remember the name.

Interviewer: Wild Wild Country?

Interviewee: No — I think so, actually — it was, like, famous three years ago, where they bought, like, an entire city — yeah, this land — and then they built their thing, and then they were, like — yes.

---


\subsection{Interview b1-aug08-2137}\label{int:b1-aug08-2137}
\subsubsection*{Metadata}
\begin{itemize}
  \item \textbf{ID}: b1-aug08-2137
  \item \textbf{Date/time}: 2026-08-08T21:37
  \item \textbf{Method}: in person
  \item \textbf{Language}: English
  \item \textbf{Location}: Germany
  \item \textbf{Interviewee}: age 35, male, Mexican nationality, main language Spanish
  \item \textbf{Questions / Answers}: 8 / 8
  \item \textbf{Total word count}: 789
  \item \textbf{Corrections log}: \texttt{cleaned/b1-aug08-2137/corrections.log}
\end{itemize}
\subsubsection*{Transcript}
INTERVIEW — Aug 8, 21:37

Interviewee: 35 years old, male, main language Spanish, speaking in English, lives in
Germany, Mexican nationality.

---

Interviewer: Perfect, all right. So, the main question is: when you hear the word cult, what comes to your mind?

Interviewee: My first thing that comes is probably the Ku Klux Klan, because — yeah, stereotypes, that way.

Interviewer: Interesting — that's what, you know, because of the movies, I'd say?

Interviewee: Yeah, yeah, but yeah — nothing more than that.

Interviewer: So the Ku Klux Klan — for you, what aspects do you think make it a cult?

Interviewee: Because it's, like, from the stereotype — particularly against, you know, a certain topic, and it's like — kind of a privilege to be in it. They sell it that way, okay. And, top — like, not everybody can be there, and, like, kind of a higher status. Yes, if you want to call it that way. And, yeah — say, like, cults — they're just like a really, really niche, small group that's against a topic. And I always think, from a bad perspective, that they are getting something that they shouldn't — being against something, that way.

Interviewer: So, if I understand well, the main feature for you is this — opposing something?

Interviewee: The way I've been shown it, like a stereotype — it's because it's like a privileged, small group of people — like, how do you call this — the one where there's only one dollar bill, you know, like this Illuminati thing, you know. It's not that it's a cult, but when I hear "cult," it's like something small, not good — I always think a cult is not the best group to belong to. No, no, okay — that's my personal opinion.

Interviewer: Yeah, yeah, of course.

Interviewee: And, and — also religious — the religious part, we probably didn't touch on that one, but a cult can also be extremely religious, in a way that it's a cult if it's a religious brand. It can also be called a cult if it's extremely into a certain religious topic, right? You know, like, these people have — I don't know — thousands of children, because it's [encouraged] by nature. And —

Interviewer: Yeah, yeah, do you have an example in mind?

Interviewee: Not that it's necessarily a cult, but — I don't mean a particular religion, but, you know, like these that are things of the church, whatever you call it, temple, whatever religion — but the cult, when they are so extreme on that religion — you know, I can also belong to that religion, but that's like a thousand kilometers apart from — exactly, that's the difference — I believe in it, and they're more extreme, in the sense of practice and ideology and all of that. I would say, okay, they also practice the ideology so strictly that they kind of raise their people to not know anything else than that, and they're also not open to the world. You know, an example — it's not a cult, a bit far from the topic, but — in Mexico there was this "La Luz del Mundo" religion, up in Guadalajara.

Interviewer: What's it called? "Luz del Mundo"? "Los del Mundo"?

Interviewee: Yeah, it's a religion that was created in Guadalajara, Mexico, and they have a church, like, I don't know, an impressive building, a lot of money there, and the people, the community — not a cult, because "community" versus "cult" — you know, like a bunch of people, and this was, of course, a broader topic. But there was an extreme part of the religion where the people who grew up there — like, if your parents were there, you grew up there, and you don't know anything else than that, so anything outside of that is wrong. And it's a bit of — okay, an extreme part of any kind of religion, whatever you believe in, you know.

Interviewer: So, then, if you were to compare this church with the Ku Klux Klan — how are they different, or how are they similar, also, because of the extreme way of practicing their beliefs?

Interviewee: Okay, no — that can be — this goes more into the religious and against a particular topic, but in a super extreme way. Yeah, no — talking about back in the day, again — this Ku Klux Klan thing is really more like a movie stereotype of what a cult means, because probably this is how I learned what the word means, you know — that can be the first thing that came to my mind when you say the word.


\subsection{Interview b1-aug08-2145}\label{int:b1-aug08-2145}
\subsubsection*{Metadata}
\begin{itemize}
  \item \textbf{ID}: b1-aug08-2145
  \item \textbf{Date/time}: 2026-08-08T21:45
  \item \textbf{Method}: in person
  \item \textbf{Language}: English
  \item \textbf{Location}: Germany
  \item \textbf{Interviewee}: age 36, male, French nationality, main language French
  \item \textbf{Questions / Answers}: 8 / 8
  \item \textbf{Total word count}: 650
  \item \textbf{Corrections log}: \texttt{cleaned/b1-aug08-2145/corrections.log}
\end{itemize}
\subsubsection*{Transcript}
INTERVIEW — Aug 8, 21:45

Interviewee: 36 years old, male, main language French, speaking in English, lives in
Germany, French nationality.

---

Interviewer: Okay, so the question is: when you hear the word cult, what comes to your mind?

Interviewee: I think the first thing that comes to my mind is more, probably, the bad side of it — the bad side of it. Because I guess, maybe it's more like growing up — you know, you were hearing these stories about girls, and the more sensational things around it — this girl, they were doing this, and it was really bad. And then there were stories, or like those people that are being set on fire because they believed in something — you know, you were told it's always crazy. And then, obviously, growing up, you learn a bit more about it, so I gradually got more interested, and then — I'm at a state now where I'm not thinking of joining a cult anytime soon, but I have more of an understanding of why you would want to believe in different things and maybe follow a different way of life. And then it also makes you question: what makes something considered a cult versus not a cult? And, you know, where do you draw the line — different societies have very different views of what would be considered a cult in certain societies versus others.

Interviewer: Okay, so I'm guessing that you already thought about the topic a bit?

Interviewee: Yeah.

Interviewer: So, when you hear the word cult, then — what example did you have in mind?

Interviewee: I think, maybe, uh — Scientology, like the [founder] — Richard... is that his name?  Scientology? No, no — the Children of God?
  No, not Children of God — sorry. The one that's displayed in the Tarantino movie.

Interviewer: Which one is that?

Interviewee: Yeah, yeah — like the Linda Kasabian story, and all that — the Manson family.

Interviewer: Yeah, Manson —

Interviewee: I said "Richard Branson" [laughs, correcting a slip].

Interviewer: Do you think Richard Branson can also be a cult leader?

Interviewee: Probably, anyway.

Interviewer:  Um, yeah — so, you thought about Scientology, and the Manson family. So, for example, with these groups — what do you think makes them a cult?

Interviewee: Well, I think they've been labeled as a cult because, I guess, it's been somehow proven that the main scheme was to have a charismatic leader who was worshipped — in a kind of messed-up way, maybe, okay — and also, a lot of it was about, you know, getting money from the members, okay, which I guess is more the bad aspect of it. I've learned also —  my favorite artist, actually, he's a guy that grew up in a cult, so I've got, like, lots of stories also from him. He really grew up in the Children of God, okay — in Europe — and managed to escape. And there's also this weird story about him, because the cult didn't allow access to doctors and medicine, and one of his younger brothers actually died when he was young. So, yeah, kind of those stories of people trying to escape those cults.

Interviewer: And can you say more about the Children of God — like, what does he say about it?

Interviewee: Um, it's a religious-adjacent group, so they have certain beliefs in common with the Bible, but probably a different interpretation of it. And I think they go around and sing a lot of Christian songs, trying to evangelize, but from a certain angle of Christianity. And, yeah, so they are very isolated, in a way — they need to follow strict principles, there are a lot of rules, and very little contact is allowed outside of the group. Yeah, so, yeah, a very different way of getting raised, I guess.


\subsection{Interview b1-aug09-2048}\label{int:b1-aug09-2048}
\subsubsection*{Metadata}
\begin{itemize}
  \item \textbf{ID}: b1-aug09-2048
  \item \textbf{Date/time}: 2026-08-09T20:48
  \item \textbf{Method}: in person
  \item \textbf{Language}: English
  \item \textbf{Location}: Germany
  \item \textbf{Interviewee}: age 34, male, Ukrainian nationality, main language Russian
  \item \textbf{Questions / Answers}: 18 / 20
  \item \textbf{Total word count}: 1157
  \item \textbf{Corrections log}: \texttt{cleaned/b1-aug09-2048/corrections.log}
  \item \textbf{Notes}: Contains profanity and explicit content preserved verbatim per instructions (discussion of the Russian Orthodox Church and MK Ultra-adjacent topics).
\end{itemize}
\subsubsection*{Transcript}
INTERVIEW — Aug 9, 20:48

Interviewee: 34 years old, male, main language Russian, speaking in English, lives in Germany, Ukrainian nationality.

---


Interviewer: Question one: when you hear the word cult, what comes to your mind?

Interviewee: All those smaller church denominations with a charismatic leader and their own little special lore — something that usually exploits vulnerable people, gullible people, you know — extracts resources and sex from them, depending on the type of cult.

Interviewer: Okay, yeah, very good. Do you have a specific example in mind — like, the first one that would come to your mind?

Interviewee: I mean, honestly, all of them do that — even the ones that are religions do it, okay. So, you know — it's like this: the difference between cult and religion, in my opinion, is as superficial as the difference between dialects and languages. You know — the ones  have armies are called languages [nodding to the saying "a language is a dialect with an army and a navy"]. So I think it's the same: as soon as the cult becomes very, very big, it's no longer called a cult, but more of a religion. But the nature is always the same. So a religion would be a cult for you, kind of, because they also have, at their inception, charismatic leaders — then they die and become saints, or prophets, or Jesuses, or Muhammads, and so forth, or Buddhas. But it always starts with some dude — there's always a dude somewhere, or a lady, in some other ones, okay. But, yeah, it's a system of belief that spreads and exploits the gullible.

Interviewer: So, like — you've mentioned a lot of different — yeah, I'm sorry —

Interviewee: No, no, it's exactly what I need, right?

Interviewer: You've mentioned a lot of different attributes — do you have in mind a specific group, or something specific, that would fill all of these attributes you mentioned? Like — you want an example?

Interviewee: Yeah, like — Scientology, the Church of Scientology, yeah. Like Mormons, yeah. Or, like, this guy — Osho, or whatever the fuck he has. That's probably it. So, I mean, there's a lot of them — like, I'm listening to podcasts about them sometimes, but I don't know their names, because they're also not memorable, you know, right? But, you know, like — the Westboro Baptist Church is one I've seen a show about.

Interviewer: what do they do —

Interviewee: something nasty, something really nasty.

Interviewee: But yeah, like — there is, nominally, some kind of difference between religions and cults, but I think it's a superficial difference, and it's like — you know, just size-based, and length of time — you perpetuate yourself in history — but there's no fundamental difference, I think. Yeah, it's just — every cult is a future religion, unless it dies out before it has a chance to become a religion. Yeah, yeah.

Interviewee: I know Jehovah's Witnesses — I mean, their religion now used to be a cult.

Interviewer: And because, like, all of the examples you gave me — I mean, you started by, I think, immediately putting it close to religion. So, like, when we start talking about the cult, it's kind of the immediate link you make?

Interviewee: Yeah, yeah, yeah, yeah, yeah — for me, it's, uh, it's kind of — every cult is a larva of a future religion-butterfly. Yeah, I think you cannot skip to religion without it first becoming a cult, okay. I think these twelve fucking disciples walking around and speaking in tongues were culty at the beginning, and now, nobody — everybody's just taking this shit for granted these days. Okay, there were these twelve weird dudes, they just walked the earth and spoke in tongues — culty as shit, in my opinion.

Interviewer: So it's mostly, like, you have a negative association with.

Interviewee: Yes

Interviewer: Absolutely, okay, have you ever encountered a cult in your life?

Interviewee: Well, I've been a member of a religion.

Interviewer: Oh, really?

Interviewee: Yeah, yeah, I'm actually still a subdeacon of the Russian Orthodox Church.

Interviewer: Really?

Interviewee: Yeah, I mean, I'm not, like, a practicing guy — yeah, since I was 19 or so — but they never expelled me. So, technically, I'm still the guy who's allowed to, you know, wipe the blood of Christ from your lips after you consume it.

Interviewer: Out of all these things you mentioned about cults in general — like, how many — which ones would apply to this, like, the Orthodox Church that you were part of?

Interviewee: All of it.

Interviewer: All of it?

Interviewee: Absolutely, all of it.

Interviewer: So which ones — you said central leader?

Interviewee: Check.

Interviewer: exploitation?

Interviewee: Check.

Interviewer: getting money, and sex?

Interviewee: Check, check.

Interviewer: Please check.

Interviewee: Well, I mean, it's not — the Russian Orthodox Church is not as known for sex scandals as, let's say, the Catholic Church, but it doesn't mean they're not there. Like, we have next to no clue what the f*** is going on in those monasteries. Yeah, but I would be shocked to discover that — whenever there's a power imbalance, people start touching each other in inappropriate places. And this is a strictly hierarchical system, so, of course, there's going to be power imbalances, and whenever there's a power imbalance, there's going to be exploitation. Yeah, it can be just your boss at work, it can be your military commander, it can be your church commander — your cult commander — but as long as you make people unequal, and somebody's going to be a little afraid to tell you when you touch me on my nipples — you might eventually touch me on my nipples. Yeah, yeah, yeah, that's my opinion.

Interviewer: Yeah, you understand —

Interviewee: And I just — I cannot give specific links to specific cases of this kind of behavior in the Russian Orthodox Church, but I have certainly heard of this in life in general.

Interviewer: Yeah, yeah, of course.

Interviewee: Yeah, there was one known subdeacon — no, no, not subdeacon — deacon, Andrei Kuraev, that's been expelled from the Russian Orthodox Church for actually being a pain in their ass. He raised this subject specifically — that we should stop pointing fingers at Catholics for being pedophiles and whatever, because we have our own hands kind of dirty too. And then they expelled him — it's really smart, dude, yeah, crazy. And he left the church too — well, he wasn't actually kicked out of the church for this specifically, but I think he was mostly not on board with, like, all the invasions of other countries that the Russian Orthodox Church otherwise stands by. Yeah, so I think that's the main reason they divorced him, not because of the sex stuff.



\subsection{Interview b1-aug09-2124}\label{int:b1-aug09-2124}
\subsubsection*{Metadata}
\begin{itemize}
  \item \textbf{ID}: b1-aug09-2124
  \item \textbf{Date/time}: 2026-08-09T21:24
  \item \textbf{Method}: in person
  \item \textbf{Language}: English
  \item \textbf{Location}: Germany
  \item \textbf{Interviewee}: age 32, male, American / German-American (see notes — no notes section present in source) nationality, main language English
  \item \textbf{Questions / Answers}: 1 / 3
  \item \textbf{Total word count}: 230
  \item \textbf{Corrections log}: \texttt{cleaned/b1-aug09-2124/corrections.log}
  \item \textbf{Notes}: Nationality field says '(see notes)' in the source but no notes section exists there — left verbatim, flagged as a data-quality gap. Interview is short (one topic: bike polo as a 'cult').
\end{itemize}
\subsubsection*{Transcript}
INTERVIEW — Aug 9, 21:24

Interviewee: 32 years old, male, main language English, speaking in English, lives in Germany, American / German-American nationality (see notes).


Interviewer: Um, when you hear the word cult, what comes to your mind?

Interviewee: Bike polo, okay — say, the specific club I currently belong to. Yeah, I guess it's culty, because it's very — it's very niche. It's a random... — I don't know, it's a subset, sub-world, and subset of society, in a way. I guess it's not really a belief system, except that we have the "shuffle gods" [referring to the random in which teams are selected in bike polo].

Interviewee: And it's a very close-knit community with little chapters around the world. We're all encouraged to go to other locations to meet up with other members of the community periodically, and come together and frolic in our celebration of the world, and life, and our existence.

Interviewee: And, yeah, it's interesting — we love the same traditions and rituals everywhere in the world. It's always interesting to realize, like, how did it get there? How do we all do the same thing? Do we all shuffle the same way? We have a lot of the same — what's it called, the term in German... "Selbstverständlichkeit" ? — concepts that we kind of take for granted, yeah. It's culty in that way.



\subsection{Interview b1-aug11-2011}\label{int:b1-aug11-2011}
\subsubsection*{Metadata}
\begin{itemize}
  \item \textbf{ID}: b1-aug11-2011
  \item \textbf{Date/time}: 2026-08-11T20:11
  \item \textbf{Method}: in person
  \item \textbf{Language}: English
  \item \textbf{Location}: Germany
  \item \textbf{Interviewee}: age 32, female, Italian nationality, main language Italian
  \item \textbf{Questions / Answers}: 4 / 9
  \item \textbf{Total word count}: 938
  \item \textbf{Corrections log}: \texttt{cleaned/b1-aug11-2011/corrections.log}
\end{itemize}
\subsubsection*{Transcript}
INTERVIEW — Aug 11, 20:11

Interviewee: 32 years old, female, main language Italian, speaking in English, lives in Germany, Italian nationality.

---

Interviewer: Well, there's basically just one main question, and the question is: when you hear the word cult, what comes to your mind?

Interviewee: To my mind comes the image of, kind of, a group of people that decide to choose to join, let's say, a community based on dogmatic beliefs. And, to me, the concept of cult also comes with the idea of a lack of freedom

Interviewee: so I have the feeling that when someone decides to join a cult, they might not even call it a cult, actually — because I also feel the name itself comes with a negative connotation. So I don't know if people who join actually say "they belong to a cult" or something — it's more like a word given from outside, by outsiders, to address this group.

Interviewee: So it's a group of people based on very specific beliefs, maybe goals, maybe a project — maybe there's a common project that people decide to commit to — but it comes with, like, compromising: once you're in, you leave something outside, you have to give up something else to be part of that community that's then called a cult. But I have more the feeling it's a name that's given to this group from outsiders.

Interviewer: And would you have specific examples of groups that would fit this definition you gave?

Interviewee: The first thing that comes to my mind, if I think about — Jehovah's Witnesses, yes, to me, I would say that sounds a bit like a "culty" community group. And also, the things that make me address them as such is the fact that they keep trying to include more and more people, and they exert pressure on people to make them join their group — which, I take for granted, comes also with a financial commitment, for example, I don't know.

Interviewee: Cult, for me, also means a lack of freedom. It means committing to rules, committing to maybe a financial engagement, and it comes with giving up some individual freedoms, let's say.

Interviewer: Do you feel like the Jehovah's Witnesses verify all the criteria that you gave, or do you have other examples — and if not, do you have other examples of a group that would fit these criteria?

Interviewee: No, I don't have much personal experience. I think it's Jehovah's Witnesses, because I've happened to interact with them. I also happened to interact with another — like, some other girls proselytizing, they stopped me on the street once, in Mitte, and they were, I think, kind of Mormons — so, a U.S.-based Christian community, "the New Jesus something", I don't know — but they looked like they were, like, 20 years old and 18, and they were sent on a mission to spread the word. And, for me, this sounds like a cult, because they were way too young to have a concrete, individual opinion about such a big topic as spirituality can be. And so, for me, this sounds like a bit of brainwashing that's been done at an age where there wasn't even yet self-consciousness and self-awareness — and the fact that they were going around, having come from another state far away to spread the word, it didn't feel like it was the result of their own decision, a choice, or the result of research — because, when it comes to spirituality, it should be the result of an individual's own search.

Interviewee: So — I'm not sure of the name of this community, maybe "the New Garden," "the New Jesus," something — I think they were Mormons, they were dressed like members of the Church of Jesus Christ of Latter-day Saints, something like this, yeah. This comes to my mind. I also met other people coming from more, like — what we can call, like, "new age" kinds of things — so, like, '60s, '70s kind of spiritual movements. So there is always, like, a connection with spirituality in what "cult" means to me, because often those communities, as far as I know, develop around — I mean, spirituality, let's say. It also happened, we had to talk with someone who was there at the beginning, in the '60s-'70s, of this Human Design kind of thing — yeah, talking to this person, it felt like — you know, for me, it's not about having strong beliefs and following them out of personal resonance, it's the fact of trying to convince others when there's no question, when there's an explicit demand. This, to me, starts to sound like a cult — and it starts to sound like, they want me in also because they need something from me, which may turn out to be money.

Interviewer: Do you think a cult needs this aspect to be a cult? Can you be a cult without trying to get new people in?

Interviewee: Probably, yes — because "cult" comes, it goes together with this perception of cult meaning also a lack of personal freedom that you give up by joining something, at a price — in exchange, they give you something greater, but at a price.

Interviewee:  For me, "cult" has a negative connotation. but I might be wrong. Society for its own reasons, has been selling the term "cult" with a negative connotation, you know. I'm quite sure — and I never did research on this — I'm quite sure that "cult" itself doesn't have a negative connotation.

---


\subsection{Interview b1-aug12-1231}\label{int:b1-aug12-1231}
\subsubsection*{Metadata}
\begin{itemize}
  \item \textbf{ID}: b1-aug12-1231
  \item \textbf{Date/time}: 2026-08-12T12:31
  \item \textbf{Method}: in person
  \item \textbf{Language}: English
  \item \textbf{Location}: Germany
  \item \textbf{Interviewee}: age 53, male, German nationality, main language German
  \item \textbf{Questions / Answers}: 11 / 12
  \item \textbf{Total word count}: 633
  \item \textbf{Corrections log}: \texttt{cleaned/b1-aug12-1231/corrections.log}
\end{itemize}
\subsubsection*{Transcript}
INTERVIEW — Aug 12, starting 12:31

Interviewee: Gender Male, Age 53, German nationality, main language German, speaking in English, lives in Germany.

--

Interviewer: So, first and only question — when you hear the word "cult," what comes to your mind?

Interviewee: Honestly, some kind of new wave band — Blue Öyster Cult, okay, this kind of stuff. And then secondly, religious cults.

Interviewer: Okay, do you think there is an actual cult aspect to Blue Öyster Cult?

Interviewee: I don't know.

Interviewer: Okay, no?

Interviewee: It's also quite different if you think about the German "Kult" or the English "cult" — because, there was this whole thing, "oh, this is cult, tick or no" [unclear phrase]. I mean it more as a noun than an adjective, like in English. I'm a non-native English speaker, so [my answer] might be distorted.

Interviewer: Yeah, no, that's fine. And then, you were talking about religious cults — do you have a specific example?

Interviewee: Yeah, this kind of mass-suicide.

Interviewer: Oh yeah, from the '70s or '80s?

Interviewee: That kind of thing, or — the Hare Krishnas.

Interviewer: And what is, in your opinion, "culty" about them — like, what are the attributes that make them a cult?

Interviewee: Well, usually following a leadership that tells you what's right and what's wrong, okay. And then identification with the leadership figure and the cult itself, so you clearly know that you're a member and that others are not.

Interviewer: Yeah.

Interviewee: And and I guess, suggest that it's a bit touching on the metaphysical — whereas the cult is more touching on the subcultural, or something like that.

Interviewer: Uh-huh. Okay — so you're making a difference between cult and religion?

Interviewee: Yeah, at first thought I would differentiate cult and religion, but then on second thought — of course, all the established religions are also cults. They just managed to be perceived a bit better.

Interviewer: Okay, so there's an aspect of cult that would be, um, like this kind of legitimacy, or how it's perceived as a whole?

Interviewee: It is — again, it's a bit difficult between German and English, I think, but for cults in Germany I would use the word "Sekte" [sect].

Interviewer: Yeah.

Interviewee: "Sekte" is kind of the differentiation from a proper religion — that it's a bit more fringe. In Germany it's quite clear: there are some churches which are registered with the state, they have the "Kirchensteuer" and so on, and once you're outside of that, you're very quickly perceived as a "Sekte," which is already connotated negatively. Nobody talks about a "Sekte" as positive, and I think that could be kind of the same as what English people mean when they say "cult." That's where it gets a bit unsharp, because I'm not a native speaker. So in German, if you say "Kult," you mean something that could also be more pop-cultural, from a movie, or — there's a cult around certain bands, but you would not say that about a sect, right.

Interviewer: And when you think about the German word "Sekte," do you have some specific ones in mind?

Interviewee: Yeah, I think — "Zeugen Jehovas", even though they're not... they're kind of also a proper church, but they don't have the same recognition. The classic would be the Hare Krishnas, Bhagwan movement, yeah, these things that kind of go viral for a certain time, and then you see them.

Interviewee: You clearly have the Catholic Church and the Protestant Church, and then the Jews also have that status. The Muslims don't properly get that status, and I think it's also because they don't have a central register. The Catholic Church is very clear, it's very hierarchical — you always know who's responsible.

---




\subsection{Interview b1-aug13-1342}\label{int:b1-aug13-1342}
\subsubsection*{Metadata}
\begin{itemize}
  \item \textbf{ID}: b1-aug13-1342
  \item \textbf{Date/time}: 2026-08-13T13:42
  \item \textbf{Method}: in person
  \item \textbf{Language}: English
  \item \textbf{Location}: Berlin, Germany
  \item \textbf{Interviewee}: age 35, female, Syrian nationality, main language Arabic
  \item \textbf{Questions / Answers}: 10 / 11
  \item \textbf{Total word count}: 1554
  \item \textbf{Corrections log}: \texttt{cleaned/b1-aug13-1342/corrections.log}
  \item \textbf{Notes}: Contains profanity preserved verbatim per instructions.
\end{itemize}
\subsubsection*{Transcript}
INTERVIEW — Aug 13, starting 13:42

Interviewee: 35 years old, female, main language Arabic, speaking in English, lives in
Berlin, Germany, Syrian nationality.

---
Interviewer: Perfect. Okay, so, first and main question — when you hear the word cult, what comes to your mind?

Interviewee: Sometimes it comes to me that, you know, kind of privileged people — that they have the privilege to, geographically, choose what to believe in. Because I believe that some places in the world, especially — if I'm talking from the perspective of the Sunna aka the Middle East — you have the religion aspect, whatever it is; I'm not talking about Muslims, even Christians or Jews — the beliefs of the regional religion are dominant over your own beliefs. But when you go more to the West, especially Western culture, you realize that people tend to choose an alternative belief to follow. Yeah, but in general, when we say "cult," I think it's a privilege to choose to be in a cult.

Interviewer: Do you have specific examples, then, of cults — I mean, given this notion of privilege?

Interviewee: For instance, I watched once a documentary about a cult called Raëlism. Have you heard of it? It's pretty crazy, nonsense shit, to be honest. But, you know, here's where I see the privilege — I believe, in some areas of the world you are dominated by the conditions around you, the daily struggles in life, where people have to be in fight-or-flight mode, so they don't have the privilege to think or rethink their beliefs and say "oh, this doesn't fit me, I'd like to follow this cult or that belief" — while in some areas of the world, especially in war zones, when people are very close to death, I believe they'd rather stick to what they were raised on, because they believe this is the remedy, or this is what gives them the better life they'll have after. So, from that point of view, I see it as a privilege — when people are safe, when they have the freedom of choice and opinion and all of that, they have the privilege to say "no, I want to be out of this religious system, and I would like to follow an alternative belief, or cult, or whatsoever." And that doesn't necessarily mean the cult is the better choice — it can be — but at least you have the privilege, the freedom, to say "I want to go out of a certain system."

Interviewee: I watched once a documentary on Netflix called "How to Become a Cult Leader," because — it's like a joke I make with my friends — that in my retirement I would like to create a cult for me and my friends, that believes in nothing, and we'd have this space where we believe in nothing together, you know. But it's very interesting how people can adopt, or follow, a certain — I don't know — ideology, just because they can, you know.

Interviewer: Yeah, so you're using both the words "ideology" and "belief" — so, do you feel like cults can be tied to either, or both?

Interviewee: Not necessarily, but it can be — it can be both, and it can be neither. But most of the time, to convince people to be part of something, you need a certain ideology that resonates with people on a certain level. Like, if we're talking about Raëlism, for instance — for people who were feeling alienated, or not fitting in — that's also very interesting about cults: that they look at the weaknesses in people's psychology and try to use that to their benefit. Like for Raëlism, where they started convincing people, "if you're feeling alienated, that's not out of nowhere — that's because we are actually created by aliens," you know? So, yeah, psychologically, it's very interesting how cults work — that they try to find the weakest spot in your brain and play around it to convince you that, no, it's something special.

Interviewer: Yeah, yeah. Yeah, cool — it's cool, I think that's enough. Yeah, thank you.

Interviewee: I'm just going to talk a little bit about the different sects of Muslims. For instance, Sunnis are prohibited from drinking alcohol, but the Alawites, for instance, don't see that as prohibited — they drink alcohol. And this is where I see the difference between a sect and a cult: sects still fall under a higher umbrella of rules, or a belief that says "this is the core belief," and the way you practice that belief can vary from one sect to another. Whereas in a cult, you don't need that — you can create the umbrella you want.

Interviewer: But — can a sect be a cult?

Interviewee: I mean, theoretically, yes, I think — or a sect would be a cult under a religious umbrella, probably. But then, I don't know what makes a cult — for me, in my understanding, a cult isn't necessarily linked to a religion, but a sect is always tied to fall under a certain religion.

Interviewer: Yeah.

Interviewee: But, for instance, there's also the Druze in the Palestine area, which we still don't know whether to call a sect or a cult, because the Druze have their own beliefs, but at the same time, one of their rules is that — depending on where they are geographically — they blend with the majority religion there. So, for example, the Druze in Syria — they're in the south of Syria, and the majority there are Muslim, so the Druze consider themselves Muslim in that context. While the Druze in Lebanon — they're in an area where the majority is Christian, so they practice and do some religious acts as Christians. But they would still call themselves Druze, not Christian, and not Muslim. So Druze, I think, could be more a cult than a sect. But Catholic or Shia is still a sect, under Muslims or Christians.

Interviewer: Yeah. Yeah.

Interviewee: I was — I don't know, like the Freemasons, actually — but still, I would see them as a cult, you know, because there's this kind of rules and regulations to be part of it, but you don't necessarily have to believe in one religion — you can be Muslim and a part of the Masons, but you can also be Christian and part of the Masons.

Interviewer: So, on which aspect, then, are the Masons a cult?

Interviewee: I'm not — I don't have a lot of knowledge about what their mission is, or what their core goal is. There was this... what was it called? What was the name of this guy who got in a scandal, for scamming people — Gomez. Gomez, the guy who loved doing Botox. Yeah, this guy [showing a picture on her phone].

Interviewer: So, can you tell me more?

Interviewee: Uh, like — ah, the Buddhafield cult — it was also about, you know, trying to bring spirituality, meditation, and enlightenment into, you know, reaching a higher spiritual space, or place. But actually, what was the scandal about, with this guy: he had, like, a really large following — I don't know, 100,000 followers — really, a lot of people, and he made this cult, even with Hollywood celebrities. And in his practices, one of the rituals they did — they sat in a dark room and meditated — and they were interviewing one of his ex-followers, and he was saying that during the meditation, Gomez would go around the room and put his thumb between people's eyes, and then people would start to see light. And then, this guy — once, he opened his eyes during the ceremony, and he realized that Gomez had an LED light on his tongue that he could turn on immediately when he touched them. So it was crazy how people who didn't know that would totally buy it, you know, and blindly believe it — it's just like the placebo effect: when you convince people so hard that this is reality, they take it. And I don't know if it comes from a highly convincing people, or because the human mind is just so fragile to be influenced.

Interviewer: So this one in particular, you'd say is a cult because of this influence — this lie, trying to influence people?

Interviewee: Yeah, and I think what makes the difference between cult and religion — I mean, I'm not a religious person, so I also can't be very sure whether prophets were there or not — but at least there is a historical record about religions. But with cults, it's not — it's just a person who wakes up one day and says "hey, I found this," and builds a whole movement around it, you know. But — as see on Wikipedia, it says it's a religious movement. Yeah, "Holy Hell" [the name of the documentary]. Yeah, but there was also a scandal about this guy, that he was using his followers because he was a Botox and plastic-surgery addict — followers, from the girls, of course — he used to convince them to try the surgeries on their bodies. So there's a lot of brainwashing and abuse allegations.


\subsection{Interview b2-aug13-1832}\label{int:b2-aug13-1832}
\subsubsection*{Metadata}
\begin{itemize}
  \item \textbf{ID}: b2-aug13-1832
  \item \textbf{Date/time}: 2026-08-13T18:32
  \item \textbf{Method}: in person
  \item \textbf{Language}: English
  \item \textbf{Location}: Germany
  \item \textbf{Interviewee}: age 44, male, Canadian nationality, main language English
  \item \textbf{Questions / Answers}: 8 / 11
  \item \textbf{Total word count}: 710
  \item \textbf{Corrections log}: \texttt{cleaned/b2-aug13-1832/corrections.log}
  \item \textbf{Notes}: File's own TRANSCRIPTION/EDITING NOTES reference names ('Mr. Mooney', 'Annie') absent from the current dialogue text — likely stale relative to this revision; left as-is, flagged for analysis.md.
\end{itemize}
\subsubsection*{Transcript}
INTERVIEW — Aug 13, 18:32

Interviewee: 44 years old, male, main language English, speaking in English, lives in Germany, Canadian nationality.

Interviewer: So that's the first, and only, question — when you hear the word cult, what comes to your mind?

Interviewee: I think it changes depending on what I've been doing in the last days.

Interviewer: Okay, like?

Interviewee: I would say one thing that comes to my mind, because over the last six months I've been quite a lot, um, searching for bicycle parts — secondhand bicycle parts. And it's very often that when there's some kind of old bicycle part, they will call it a "cult" bicycle part.

Interviewer: Okay — like a cult mountain bike or whatever?

Interviewee: Shit, yeah — those are usually the quite expensive things, which I, of course, never buy. Yeah.

Interviewee: And then there are also, like, cult films. So cult films are usually films that were not very economically successful, but they have a small following — very devout fans — towards films, or musicians, or albums, or artists, or whatever. So, in these cases, I think you would use "cult" as an adjective.

Interviewer: Um, yeah, if you were to use it as a noun, then, like — things that are cults — what would be examples of that?

Interviewee: Well, I suppose it's related, whether it's adjective or noun. It's — it's speaking about a certain devotion to something. So, whether it's to a derailleur for a bicycle, or a Korean Christian dude — Mr. Moon — or whether it's... anything, you know?

Interviewer: So one example that you mentioned, a specific bicycle part?

Interviewee: No, I just — it's something I came across quite a lot when looking, and people trying to sell things — they, you know, of course, want to sell it. So they would say a lot of good stuff about it, and, yeah, very often, they'd use this word "cult" to, like, say this is a really cool thing, right?

Interviewer: Okay, ah — so, like, as a positive thing?

Interviewee: Actually, yes.

Interviewer: And, uh, then you were talking about the moon...

Interviewee: Oh, this was just because, uh, of the moon — yesterday there was an eclipse, a solar eclipse, and I had to photograph that. My friend [Name redacted] took a photo of me putting the welding helmet onto my other friend [Name redacted] so she could look at the sun, right?

Interviewee: And then I sent it to her — the text — and then we started talking. She started making jokes about being really ready to fight and die for the moon, and that conversation went on. And then I remembered — oh yeah, there was some kind of cult called the Moonies in the '50s, or '60s, or '70s, or something.

Interviewee: Yeah, but I couldn't remember what it was, so I just looked it up. Oh, that's what it was, right. Okay, I thought — that's why it's on my mind right now, because of the solar eclipse yesterday.

Interviewer: That's interesting. So, like, you remember it — you remember it was called a cult, but you don't really necessarily remember exactly why?

Interviewee: No, I didn't know anything about it. I just remembered that there was this cult called the Moonies... or maybe it still is, I don't know.

---

TRANSCRIPTION / EDITING NOTES
- "Mr. Mooney" almost certainly refers to Sun Myung Moon, the Korean-born founder of
  the Unification Church — whose members are popularly nicknamed "the Moonies," a
  name the interviewee independently brings up a few lines later when describing the
  cult he half-remembers from the solar-eclipse conversation. Kept as transcribed
  with a [?] flag rather than silently renamed, since the raw audio isn't available
  to confirm.
- "Annie" is a best guess for a name that's unclear in the raw text ("Annie's" —
  possibly a mis-transcribed name, or a possessive referring to someone else
  entirely). Flagged rather than assumed.
- Recurring, unambiguous ASR errors ("a cut" for "a cult") were corrected silently
  per the pattern established across the earlier batch.
- One instance of profanity was masked in the original raw transcript ("S***") and
  has been left as-is rather than guessed at.


\subsection{Interview b2-aug13-1840}\label{int:b2-aug13-1840}
\subsubsection*{Metadata}
\begin{itemize}
  \item \textbf{ID}: b2-aug13-1840
  \item \textbf{Date/time}: 2026-08-13T18:40
  \item \textbf{Method}: in person
  \item \textbf{Language}: English
  \item \textbf{Location}: Berlin, Germany
  \item \textbf{Interviewee}: age 43, male, Italian nationality, main language Italian
  \item \textbf{Questions / Answers}: 5 / 6
  \item \textbf{Total word count}: 310
  \item \textbf{Corrections log}: \texttt{cleaned/b2-aug13-1840/corrections.log}
\end{itemize}
\subsubsection*{Transcript}
INTERVIEW — Aug 13, 18:40

Interviewee: 43 years old, male, main language Italian, speaking in English, lives in
Berlin, Germany, Italian nationality.

---

Interviewer: When you hear the word cult, what comes to your mind?

Interviewee: So, "cult" — if we translate it properly in Italian, should be "la setta". So, uh, usually it's related to a collective, a group of people following some sort of belief, and related to religion — I guess it's always connected to religion, right? I guess. But non-conventional religion, so that's why they are, most of the time, concentrated in small groups. And, somehow, when I think about cult, I think about something like dark religious — like praying to the devil, or something like that.

Interviewer: Yes

Interviewee: I don't know, I guess this is more or less the idea that I have. It's really easy to get kind of convinced to join a cult, by using some kind of false induction, false information. Somehow they try to convince people to join the cult.

Interviewer: Well — do you have specific examples of cults that, like, immediately come to your mind?

Interviewee: I mean, those ones connected to the devil — I'd just say, like, Satanists.

Interviewer: Yes

Interviewee: Or maybe Scientology is also a cult, maybe? Not sure about it.

Interviewer: What are the characteristics, you know, the details — that make a group of people set as a cult?

Interviewee: I'm not sure, actually, what is actually the right definition of it, you know. But, yeah, the first thing I think of is, like, Satan — Satan is somehow.

Interviewee: I mean, they always name it this way — like "the cult of Satan," "the cult of the Archangels," I don't know, stuff like that — so it's easy to relate it to this kind of Satanism, or dark cult.

---


\subsection{Interview b3-aug16-1514}\label{int:b3-aug16-1514}
\subsubsection*{Metadata}
\begin{itemize}
  \item \textbf{ID}: b3-aug16-1514
  \item \textbf{Date/time}: 2026-08-16T15:14
  \item \textbf{Method}: in person
  \item \textbf{Language}: English
  \item \textbf{Location}: Germany
  \item \textbf{Interviewee}: age 34, female, Greek and German nationality, main language Greek
  \item \textbf{Questions / Answers}: 3 / 4
  \item \textbf{Total word count}: 238
  \item \textbf{Corrections log}: \texttt{cleaned/b3-aug16-1514/corrections.log}
\end{itemize}
\subsubsection*{Transcript}
INTERVIEW — Aug 16, 15:14

Interviewee: 34 years old, female, main language Greek, speaking in English, lives in Germany, nationality Greek and German.


---

Interviewer: When you hear the word "cult," what comes to your mind?

Interviewee: Charles Manson — and CIA.

Interviewer: Great — can you explain what you feel is cult-like about these things?

Interviewee: So, about Charles Manson, I think — because he ran the cult, I don't know how to explain it — he killed a bunch of people, thinking that he was serving some type of ulterior motive. I guess — or not, yeah — because also, there was no reason to kill them, but he did it anyway. And he had a bunch of people running after him doing that stuff.

Interviewee: ...And MK Ultra

Interviewer:  Can you say more, like what you really associate with that?

Interviewee: So MK Ultra — was a secret program, an entirely secret American intelligence program. And there's a theory that he [Manson] was one of the subjects there, and basically they tried to create personality disorders from scratch, and they'd torture people. And there's a theory — we don't know, maybe it's a conspiracy theory — that some of them, like Charles Manson, started a cult, or like, uh, Brian Jones. Yeah, always like the '60s, '70s — and I think this kind of thing in America is also related to Cold War paranoia. Yeah.



\subsection{Interview b3-aug16-1517}\label{int:b3-aug16-1517}
\subsubsection*{Metadata}
\begin{itemize}
  \item \textbf{ID}: b3-aug16-1517
  \item \textbf{Date/time}: 2026-08-16T15:17
  \item \textbf{Method}: in person
  \item \textbf{Language}: English
  \item \textbf{Location}: Germany
  \item \textbf{Interviewee}: age 36, female, Turkish nationality, main language Turkish
  \item \textbf{Questions / Answers}: 3 / 4
  \item \textbf{Total word count}: 172
  \item \textbf{Corrections log}: \texttt{cleaned/b3-aug16-1517/corrections.log}
  \item \textbf{Notes}: Source file has no 'INTERVIEW — date, time' header (unlike other Batch 3 files); date\_time above is taken from the filename.
\end{itemize}
\subsubsection*{Transcript}
Interviewee: 36 years old, female, main language Turkish, speaking in English, lives in Germany, nationality Turkish.

Interviewer: Okay. So, when I say the word "cult," what comes to your mind?

Interviewee: When you say the word cult, it comes to my mind — religious cults in Turkey.

Interviewer: Do they have specific names?

Interviewee: Yes, I think one is led by Fethullah Gülen, and the other one is led by another guy whose name I forgot.


Interviewer: What makes them a cult, in your opinion?

Interviewee: Because they have a singular leader, it's not easy to get into the group, you have to go through an integration process, and it's a very closed community.

Interviewee: Actually — there's one more that comes to mind now, called Menzil, I think, on the east side — or something like this, east side of Turkey. They have a whole town for themselves; getting in and out is controlled. And the members have to live by the rules of the group — very strict rules.

--


\subsection{Interview b3-aug18-1645}\label{int:b3-aug18-1645}
\subsubsection*{Metadata}
\begin{itemize}
  \item \textbf{ID}: b3-aug18-1645
  \item \textbf{Date/time}: 2026-08-18T16:45
  \item \textbf{Method}: in person
  \item \textbf{Language}: French, translated to English
  \item \textbf{Location}: France
  \item \textbf{Interviewee}: age 32, male, French nationality, main language French
  \item \textbf{Questions / Answers}: 4 / 7
  \item \textbf{Total word count}: 331
  \item \textbf{Corrections log}: \texttt{cleaned/b3-aug18-1645/corrections.log}
\end{itemize}
\subsubsection*{Transcript}
\begin{otherlanguage}{french}
INTERVIEW — Aug 18, 16:45

Interviewee: 32 years old, male, main language French, speaking in French, lives in France, nationality French.

---

Interviewer : Très bien. Quand tu entends le mot secte, qu'est-ce qui te vient à l'esprit ?

Interviewee : Gourou. Groupe. Un groupe qui suit une doctrine. Qui suit des règles. Avec souvent un gourou. Ouais, c'est des religions qui n'ont pas percé, quoi.

Interviewer : Et tu as quoi comme exemple ?

Interviewee : La secte de Raël. La secte... les Témoins de Jéhovah aussi... Est-ce que c'est vraiment une secte? Est-ce que c'est pas une religion, les Témoins de Jéhovah ?

Interviewee : Est-ce que, je sais pas, ceux qui peuvent défendre... — la secte de l'homéopathie?

Interviewer : Tu penses que l'homéopathie, tu pourrais dire que c'est une secte ?

Interviewee : Non, pas sûr, parce qu'il y a pas de... après je sais pas, c'est pas forcément un critère, mais ça ne profite pas à une personne.

Interviewer : Donc ce serait plus quoi, l'homéopathie, ce serait quoi ?

Interviewee : Ce serait... non, mais, c'est pas une secte. Mais après, des sectes actuelles — est-ce qu'il y a des sectes actuellement certifiées ?

Interviewee : Pour moi, une secte, déjà, c'est qu'on ne la connaît pas tant qu'elle n'est pas vraiment partie en cacahuète, quoi. On la connaît quand il y a eu des drames — genre des suicides collectifs, des trucs comme ça — parce qu'en fait, si tu sais que c'est une secte, a priori, c'est que c'est néfaste. Et donc je suppose que les instances politiques font en sorte de fermer les sectes, quoi ? C'est pour ça que tu vois, si — en tout cas les Témoins de Jéhovah, je suppose qu'il y a des gens qui ont dû travailler dessus et se dire : bon, peut-être que c'est pas assez hard pour qu'on pour qu'on leur cloue le bec.

Interviewee : J'ai pas d'autre exemple de secte que ça.

\end{otherlanguage}
\subsubsection*{English translation}
INTERVIEW — Aug 18, 16:45 (English translation of b3-aug18-1645/transcript.txt)
Translation note: produced directly by the assistant compiling this corpus, not a
certified translation. Worth spot-checking against the French original for nuance.

Interviewee: 32 years old, male, main language French, speaking in French, lives in France, nationality French.

---

Interviewer: Great. When you hear the word "sect" [secte], what comes to mind?

Interviewee: Guru. Group. A group that follows a doctrine. That follows rules. Often with a guru. Yeah, they're religions that never broke through, kind of.

Interviewer: And what would you give as an example?

Interviewee: The Raël sect. The sect... the Jehovah's Witnesses too... Is that really a sect? Isn't it a religion, the Jehovah's Witnesses?

Interviewee: I don't know, whoever can argue for... — the homeopathy sect?

Interviewer: You think homeopathy — you could say that's a sect?

Interviewee: No, not sure, because there isn't any... then again I don't know, it's not necessarily a criterion, but it doesn't benefit one person.

Interviewer: So what would homeopathy be, then?

Interviewee: It would be... no, it's not a sect. But then, current sects — are there any officially recognized sects right now?

Interviewee: For me, a sect, really, is something you don't know about until it's really gone off the rails. You find out about it once there have been tragedies — like mass suicides, things like that — because actually, if you know it's a sect, it's presumably because it's harmful. So I guess the political authorities make sure to shut down sects? That's why, you see — anyway, the Jehovah's Witnesses, I imagine there are people who must have worked on that and thought: well, maybe it's not "hard" enough for them to be shut down.

Interviewee: I don't have another example of a sect besides that.


\subsection{Interview b3-aug21-1932}\label{int:b3-aug21-1932}
\subsubsection*{Metadata}
\begin{itemize}
  \item \textbf{ID}: b3-aug21-1932
  \item \textbf{Date/time}: 2026-08-21T19:32
  \item \textbf{Method}: in person
  \item \textbf{Language}: French, translated to English
  \item \textbf{Location}: Spain
  \item \textbf{Interviewee}: age 32, male, French nationality, main language French
  \item \textbf{Questions / Answers}: 5 / 7
  \item \textbf{Total word count}: 298
  \item \textbf{Corrections log}: \texttt{cleaned/b3-aug21-1932/corrections.log}
\end{itemize}
\subsubsection*{Transcript}
\begin{otherlanguage}{french}
INTERVIEW — Aug 21, 19:32

Interviewee: 32 years old, male, main language French, speaking in
French, lives in Spain, nationality French.

---

Interviewer : Très bien. Quand tu entends le mot secte, qu'est-ce qui te vient à l'esprit ?

Interviewee : Illuminati.

Interviewer : Ah ouais, ok !

Interviewee : Je sais pas, en fait, là je pensais un peu à une conspiration, tu vois. Mais sinon je pense à un groupe.

Interviewer : Un groupe ?

Interviewee : Une communauté.

Interviewer : Une communauté qui... ?

Interviewee : ... qui se rassemble autour d'un gourou.

Interviewer : Ok. Et du coup, les Illuminati, ça t'est venu spontanément — qu'est-ce qui fait que tu le rapproches d'une secte ?

Interviewee : Bah, je sais pas, c'est le premier mot qui m'est venu à l'esprit. Je pense que c'est un peu le côté mystérieux aussi qu'on retrouve dans la culture des sectes. Tu vois, genre, c'est un peu décrit comme quelque chose... je sais pas. Maintenant, il y a des reportages Netflix sur les sectes, les gens sont un peu friands de ça — un peu autour du secret, un peu... En fait, j'ai l'impression que les gens veulent aussi un peu de mystère, tu vois. Et je pense qu'on part un peu, on s'éloigne un peu de certains aspects scientifiques, et du coup les gens veulent donner du sens aux choses. Et du coup ils sont attirés par ces choses, un peu bizarre.

Interviewee : Enfin, qui expliqueraient beaucoup de choses, en fait, sur la société moderne, parce que...ça dévie un peu l'attention des gens. En les attirant vers quelque chose.

Interviewee : C'est un peu comme les phénomènes de réseaux sociaux. En fait, tous les gens veulent du contenu un peu hors norme, tu vois ?

---

\end{otherlanguage}
\subsubsection*{English translation}
INTERVIEW — Aug 21, 19:32 (English translation of b3-aug21-1932/transcript.txt)
Translation note: produced directly by the assistant compiling this corpus, not a
certified translation. Worth spot-checking against the French original for nuance.

Interviewee: 32 years old, male, main language French, speaking in French, lives in Spain, nationality French.

---

Interviewer: Great. When you hear the word "sect," what comes to mind?

Interviewee: Illuminati.

Interviewer: Oh yeah, okay!

Interviewee: I don't know, actually, just then I was thinking a bit about a conspiracy, you know. But otherwise I think of a group.

Interviewer: A group?

Interviewee: A community.

Interviewer: A community that...?

Interviewee: ...that gathers around a guru.

Interviewer: Okay. So, the Illuminati came to you spontaneously — what makes you connect it to a sect?

Interviewee: Well, I don't know, it's the first word that came to mind. I think it's a bit the mysterious side too, that you find in sect culture. You know, like, it's kind of described as something... I don't know. Now there are Netflix documentaries about sects, people are a bit into that — a bit around the secrecy, a bit... Actually, I get the impression people also want a bit of mystery, you know. And I think we're drifting a little, moving away a bit from certain scientific aspects, and so people want to give meaning to things. And so they're drawn to these things, a bit strange.

Interviewee: Anyway, things that would actually explain a lot about modern society, because... it diverts people's attention a bit. By drawing them toward something.

Interviewee: It's a bit like social media phenomena. Actually, everyone wants slightly out-of-the-ordinary content, you know?

---


\subsection{Interview b3-aug22-2057}\label{int:b3-aug22-2057}
\subsubsection*{Metadata}
\begin{itemize}
  \item \textbf{ID}: b3-aug22-2057
  \item \textbf{Date/time}: 2026-08-22T20:57
  \item \textbf{Method}: in person
  \item \textbf{Language}: French, translated to English
  \item \textbf{Location}: France (inferred from context, not explicitly captured — see file's own notes)
  \item \textbf{Interviewee}: age 32, male, French nationality, main language French
  \item \textbf{Questions / Answers}: 3 / 3
  \item \textbf{Total word count}: 412
  \item \textbf{Corrections log}: \texttt{cleaned/b3-aug22-2057/corrections.log}
  \item \textbf{Notes}: Interview ends abruptly with no closing exchange; treated as complete per the batch's own merging analysis (no continuation file exists).
\end{itemize}
\subsubsection*{Transcript}
\begin{otherlanguage}{french}
INTERVIEW — Aug 22, 20:57

Interviewee: 32 years old, male, main language French, speaking in French, lives in France, French nationality.


---

Interviewer : Quand tu entends le mot secte, qu'est-ce qui te vient à l'esprit ?

Interviewee : Manipulation, groupe, idéologie.

Interviewer : Est-ce qu'il y a des exemples de groupes, du coup, comme ça, qui te viendraient à l'esprit ?

Interviewee : Le truc du cercle solaire [référence à l'Ordre du Temple Solaire] ? Non, j'ai pas les noms, je sais pas trop. Si, il y en a — un peu moins sectaire, mais il y a celle avec Tom Cruise — la Scientologie. Ce qui me vient à l'esprit, c'est la Scientologie. Même les... comment on dit... putain.. New Wave? La croyance New Wave ?

Interviewer : New Age ?

Interviewee : Voilà !

---

TRANSCRIPTION / EDITING NOTES
- Nationality: the interviewer's question is present in the raw text, but the transcript jumps straight to the main question with no answer captured — genuinely missing. Country of residence is not explicitly stated either; "France" is inferred from context (the interviewer's aside reads roughly as confirming "France" as an answer) and flagged as an inference, not a captured fact.
- The "C'est pour ton master ?" / "Oui, c'est pour mon master, pas d'autres questions ?" exchange is a reconstruction of a garbled meta-aside in the raw text ("mais c'est pourquoi juste pour ton master mais c'est quoi le pose? Pas de questions?"); the gist (interviewee asking whether this is for the interviewer's Master's thesis) seems right, but the exact wording is uncertain.
- "cercle solaire" → likely refers to "l'Ordre du Temple Solaire" (the Solar Temple), a real cult responsible for a series of mass suicides/murders in Switzerland, France, and Canada in 1994–1997 — proposed with moderate-high confidence but not silently substituted since the exact name differs from what was said.
- The raw transcript itself contains a masked word ("p*****" — asterisks in the source file), i.e. the ASR system censored or couldn't render a word the interviewee said; reproduced as-is since the source is illegible/censored, not a transcription choice.
- The interview ends on "Voilà!" right after the interviewee appears to search for a (possibly redacted) word — there's no closing "thank you" exchange, and it's unclear whether the recording was cut short or the interviewee simply gave up on finding the word. Treated as a complete, if abrupt, standalone interview since no continuation file exists.

\end{otherlanguage}
\subsubsection*{English translation}
INTERVIEW — Aug 22, 20:57 (English translation of b3-aug22-2057/transcript.txt)
Translation note: produced directly by the assistant compiling this corpus, not a
certified translation. Worth spot-checking against the French original for nuance.
Only the interview dialogue is translated; the source file's own
TRANSCRIPTION / EDITING NOTES are not duplicated here (see the original transcript).

Interviewee: 32 years old, male, main language French, speaking in French, lives in France, French nationality.

---

Interviewer: When you hear the word "sect," what comes to mind?

Interviewee: Manipulation, group, ideology.

Interviewer: Are there any examples of groups, then, that would come to mind?

Interviewee: The solar circle thing [reference to the Order of the Solar Temple]? No, I don't have the names, I don't really know. Actually, yes, there is one — a bit less sect-like, but there's the one with Tom Cruise — Scientology. What comes to mind is Scientology. Even the... how do you say it... damn it... New Wave? The New Wave belief?

Interviewer: New Age?

Interviewee: That's it!

---


\subsection{Interview b3-aug22-2101}\label{int:b3-aug22-2101}
\subsubsection*{Metadata}
\begin{itemize}
  \item \textbf{ID}: b3-aug22-2101
  \item \textbf{Date/time}: 2026-08-22T21:01
  \item \textbf{Method}: in person
  \item \textbf{Language}: French, translated to English
  \item \textbf{Location}: France
  \item \textbf{Interviewee}: age 32, female, French nationality, main language French
  \item \textbf{Questions / Answers}: 4 / 5
  \item \textbf{Total word count}: 475
  \item \textbf{Corrections log}: \texttt{cleaned/b3-aug22-2101/corrections.log}
\end{itemize}
\subsubsection*{Transcript}
\begin{otherlanguage}{french}
INTERVIEW — Aug 22, 21:01

Interviewee: 32 years old, female, main language French, speaking in French, lives in France, nationality French.

---

Interviewer : Quand tu entends le mot secte, qu'est-ce qui te vient à l'esprit ?

Interviewee : Qu'est-ce qui me vient à l'esprit... des groupes fermés qui endoctrinent des gens en échange d'argent, la plupart du temps, et qui peuvent conduire à des dérives dangereuses.

Interviewer : Ok, très bien. Et est-ce que tu as des exemples, du coup ?

Interviewee : La Scientologie, le truc de Raël... Non, à part ça, pas d'autres exemples. Ah si! J'ai entendu récemment un truc — un groupe dans le sud de la France. Ouais, merde, comment ça s'appelait déjà ? Qui font, comme  dans la série là, quand on regardait... où ils font des chansons, et c'est un truc genre catho plus plus, droite plus plus,je sais plus comment ça s'appelle. les brigandes?

Interviewer : Tu as d'autres groupes, ou tu as d'autres choses à dire là-dessus ?

Interviewee : C'est pas forcément très bien.

Interviewer : Très bien. Et du coup, pourquoi ce groupe-là en particulier, tu dirais que c'est une secte ?

Interviewee : En tout cas, dans le petit reportage que j'avais vu, ça avait l'air d'être quand même une idéologie martelée très fort, et des histoires d'argent qui circulaient en faveur d'une espèce de gourou — entre guillemets, un personnage — qui était au centre de l'attention et à l'origine de la création du groupe.

Interviewee : A priori, dans la définition même, ça implique un truc négatif.


---

TRANSCRIPTION / EDITING NOTES
- The group the interviewee is trying to name (raw: "Plus plus et droit de plus plus là je sais plus comment s'appelle les briques. Ce film?") is badly garbled and could not be reliably reconstructed word-for-word — left flagged with [?] rather than filled in.
- That said, the surrounding description — a group in the south of France, seen in a TV series/documentary, associated with singing, with restricted ability to leave — matches the real Twelve Tribes community (also known in France as "Tabitha's Place"), which has a base in Toulouse and has been the subject of French-language documentaries about members' inability to leave and children's treatment. Proposed with moderate confidence as the group being described, based on a web search, but not inserted into the dialogue itself since the interviewee's actual words for the name weren't recoverable.
- "des rives dangereuses" → "des dérives dangereuses," corrected silently (clear ASR mishearing, "dangerous drifts/excesses").
- "Première et. Euh sale question." was reconstructed as the interviewer reassuring the interviewee that the (only) study question isn't a difficult/unpleasant one; the exact original wording is uncertain and flagged with [?].
- "En français?" (raw, for country of residence) → rendered as "En France," corrected silently as an unambiguous ASR slip.

\end{otherlanguage}
\subsubsection*{English translation}
INTERVIEW — Aug 22, 21:01 (English translation of b3-aug22-2101/transcript.txt)
Translation note: produced directly by the assistant compiling this corpus, not a
certified translation. Worth spot-checking against the French original for nuance.
Only the interview dialogue is translated; the source file's own
TRANSCRIPTION / EDITING NOTES are not duplicated here (see the original transcript).

Interviewee: 32 years old, female, main language French, speaking in French, lives in France, nationality French.

---

Interviewer: When you hear the word "sect," what comes to mind?

Interviewee: What comes to mind... closed groups that indoctrinate people in exchange for money, most of the time, and that can lead to dangerous excesses.

Interviewer: Okay, very good. And do you have any examples, then?

Interviewee: Scientology, the Raël thing... No, apart from that, no other examples. Oh wait, yes! I heard about something recently — a group in the south of France. Yeah, damn, what was it called again? They do, like in that series, when we were watching... where they sing songs, and it's like a super-Catholic thing, super-right-wing, I don't know what it's called anymore. Les Brigandes?

Interviewer: Do you have other groups, or other things to say about this?

Interviewee: It's not necessarily very good.

Interviewer: Very well. So, why would you say this particular group is a sect?

Interviewee: In any case, in the short documentary I'd seen, it seemed to be an ideology hammered in very hard, and stories of money circulating in favor of a kind of guru — in quotation marks, a figure — who was at the center of attention and at the origin of the group's creation.

Interviewee: Presumably, by definition itself, that implies something negative.

---


\subsection{Interview b3-aug23-1213}\label{int:b3-aug23-1213}
\subsubsection*{Metadata}
\begin{itemize}
  \item \textbf{ID}: b3-aug23-1213
  \item \textbf{Date/time}: 2026-08-23T12:13
  \item \textbf{Method}: in person
  \item \textbf{Language}: French, translated to English
  \item \textbf{Location}: France
  \item \textbf{Interviewee}: age 97, male, French nationality, main language French
  \item \textbf{Questions / Answers}: 5 / 7
  \item \textbf{Total word count}: 253
  \item \textbf{Corrections log}: \texttt{cleaned/b3-aug23-1213/corrections.log}
  \item \textbf{Notes}: Stated age (97) and gender are flagged as unreliable/possibly joking by the batch's own edited-transcripts README; reported here exactly as stated in the source.
\end{itemize}
\subsubsection*{Transcript}
\begin{otherlanguage}{french}
INTERVIEW — Aug 23, 12:13

Interviewee: age 97, gender male, main language French, speaking in French, lives in France, nationality French.

---

Interviewer : Bon, déjà ça, c'est tout bon. Et maintenant la question de l'étude — il y a qu'une seule question, et tu ne peux pas te tromper. Tu dis tout ce qui te passe par la tête, ok ? Quand je te dis le mot secte à quoi tu penses?

Interviewee : Sept?

Interviewer : Une secte.

Interviewee : Insecte?

Interviewer : une secte,

Interviewee : J'ai compris — une secte.

Interviewer : À quoi tu penses ?

Interviewee : Des groupements... semi-religieux, quand je... je ne comprends pas toujours [?]... ça n'a pas toujours été une bonne réputation, enfin... je sais... que... je ne peux pas bien compléter, je sais pas. C'est quelque chose qu'on évite de fréquenter, en général, parce que, dans mon style de vie...

Interviewer : Est-ce que tu as des exemples de sectes qui te viennent à l'esprit, quand je te dis "secte" ?

Interviewee : ...et maintenant, il n'y en a plus beaucoup — il y en a eu à une époque, je sais pas si c'est encore une secte maintenant. On va pas citer une secte en particulier..

Interviewee: On sait jamais bien à partir de quelle moment c'est une secte ou c'est est pas une
Interviewee: À partir de quel moment ça peut avoir des connotations plus ou moins religieuses, ou plus ou moins politiques?, ou autre chose... ça peut être large.

\end{otherlanguage}
\subsubsection*{English translation}
INTERVIEW — Aug 23, 12:13 (English translation of b3-aug23-1213/transcript.txt)
Translation note: produced directly by the assistant compiling this corpus, not a
certified translation. Worth spot-checking against the French original for nuance —
this interview in particular hinges on a mishearing pun in French ("secte" / "sept" /
"insecte") that does not carry over to English; a translator's note is added inline
rather than attempting an English pun.

Interviewee: age 97, gender male, main language French, speaking in French, lives in France, nationality French.

---

Interviewer: Okay, that's all good. Now the study question — there's only one question, and you can't get it wrong. Just say whatever comes to mind, okay? When I say the word "secte" [sect/cult], what do you think of?

Interviewee: Seven? [mishears "secte" as "sept," French for "seven" — homophone in spoken French]

Interviewer: A sect.

Interviewee: Insect? [mishears "secte" as "insecte"]

Interviewer: a sect,

Interviewee: I understood — a sect.

Interviewer: What do you think of?

Interviewee: Groups... semi-religious, when I... I don't always understand [?]... it hasn't always had a good reputation, well... I know... that... I can't quite finish the thought, I don't know. It's something people generally avoid associating with, because, in my way of life...

Interviewer: Do you have any examples of sects that come to mind, when I say "sect"?

Interviewee: ...and now, there aren't many left — there were some at one point, I don't know if it's still a sect now. Let's not name one in particular..

Interviewee: You never really know from what point something is a sect or isn't
Interviewee: From what point can it take on more or less religious connotations, or more or less political ones, or something else... it can be broad.


\subsection{Interview ig-02}\label{int:ig-02}
\subsubsection*{Metadata}
\begin{itemize}
  \item \textbf{ID}: ig-02
  \item \textbf{Date/time}: 2026-08-13 (date only — no time of day given in source)
  \item \textbf{Method}: Instagram
  \item \textbf{Language}: English
  \item \textbf{Location}: Germany
  \item \textbf{Interviewee}: age 27, female, Turkish nationality, main language Turkish
  \item \textbf{Questions / Answers}: 3 / 3
  \item \textbf{Total word count}: 224
  \item \textbf{Corrections log}: \texttt{cleaned/ig-02/corrections.log}
  \item \textbf{Notes}: Instagram DM interview, gathered after a story post; conducted via direct message per the file's own Interview Notes.
\end{itemize}
\subsubsection*{Transcript}
Age: 27, Gender: Female, Nationality: Turkish, Main language: Turkish, Language of interview: English, lives in Germany
--
Interviewer: When you hear the word cult, what comes to your mind?
Interviewee: Turkish Religious cults.
Interviewer: Amazing thank you. Briefly would you mind telling me why you associate them with a cult?
Interviewee: I have memories as a kid/teen where these religious groups would run tutoring centers, preparing students for the university exams but also inviting them to Quran readings and kinda radicalizing them, they would organize special camps, give scholarship etc. These would happen on sight but also kinda behind the doors and there was also several groups with different leaders. Very cult-ish, but also they were literally called cult in Turkish, which is the word tarikat. Now they cover a significant base of Erdoğan voters and also work and organize in different governmental bodies.
Interviewer: great thank you so much. do you remember some specific names?
Interviewee: Nakşibendi (one of the cults) , menzil (also one of the cults, like a subcult of nakşibendi), Süleymancilar( also subcult of nakşibendi), saidi nursi (cult leader of nursi cult)
---
Interview Notes: gathered on instagram after posting a story asking for people to share their thoughts on the word "cult." The interview was conducted via direct message, on Thursday, August 13th 2026. Personal information have been anonymized.


\subsection{Interview ig-03}\label{int:ig-03}
\subsubsection*{Metadata}
\begin{itemize}
  \item \textbf{ID}: ig-03
  \item \textbf{Date/time}: 2026-08-13 (date only — no time of day given in source)
  \item \textbf{Method}: Instagram
  \item \textbf{Language}: English
  \item \textbf{Location}: Germany
  \item \textbf{Interviewee}: age 29, male, Lebanese nationality, main language Arabic
  \item \textbf{Questions / Answers}: 1 / 10
  \item \textbf{Total word count}: 178
  \item \textbf{Corrections log}: \texttt{cleaned/ig-03/corrections.log}
  \item \textbf{Notes}: Instagram DM interview. One phrase ('OR import but underground art') is possibly a dropped-syllable typo for 'important' — left uncorrected pending confirmation (see corrections.log).
\end{itemize}
\subsubsection*{Transcript}
Age: 29, Gender: Male, Nationality: Lebanese, Main language: Arabic, Language of interview: English, lives in Germany
--
Interviewer: When you hear the word cult, what comes to your mind?
Interviewee: Weird people coming together to do weird things. OR import but underground art
Interviewee: Charles Manson and his group
Interviewee: Satanist cults?
Interviewee: That's what they're called no? Some people use cult/sect interchangeably
Interviewee: I wouldn't say Shi'a muslims are a cult, they are a religious sect
Interviewee: But they sometimes have cult-like activities like self-flagellation
Interviewee: Similarly I wouldn't think of catholicism as a cult, but a religious sect
Interviewee: But when catholics in Ethiopia (to give an example) practice self-flagellation or self-crucifixion then it feels more like a cult
Interviewee: And re: cult in the arts
Interviewee: Ziad Rahbani's work in theater and music is 'cult' in Lebanon
---
Interview Notes: gathered on instagram after posting a story asking for people to share their thoughts on the word "cult." The interview was conducted via direct message, on Thursday, August 13th 2026. Personal information have been anonymized.


\subsection{Interview ig-04}\label{int:ig-04}
\subsubsection*{Metadata}
\begin{itemize}
  \item \textbf{ID}: ig-04
  \item \textbf{Date/time}: 2026-08-13 (date only — no time of day given in source)
  \item \textbf{Method}: Instagram
  \item \textbf{Language}: English
  \item \textbf{Location}: Germany
  \item \textbf{Interviewee}: age 38, female, Portuguese nationality, main language Portuguese
  \item \textbf{Questions / Answers}: 3 / 7
  \item \textbf{Total word count}: 231
  \item \textbf{Corrections log}: \texttt{cleaned/ig-04/corrections.log}
  \item \textbf{Notes}: Instagram DM interview. Contains deliberate leetspeak stylization ('isr43l') preserved verbatim as the interviewee's own phrasing, not corrected.
\end{itemize}
\subsubsection*{Transcript}
Age: 38, Gender: Female, Nationality: Portuguese, Main language: Portuguese, Language of interview: English, lives in Germany
--
Interviewer: When you hear the word cult, what comes to your mind?
Interviewee: People dressed in white
Interviewer: Amazing thank you. Briefly would you mind telling me why you associate them with a cult?
Interviewee: There's often someone dressed in white, pope, orthodox ppl, imam, kkk, umbanda, evangelical, jesus underwear is white too. The sacred dove  that represents the spirit is white, usa isr43l have white, hygiene standards are linked to white. At the supermarket whitening products, medical machines are white, stores are white.
Interviewee: People marry dressed in white, purity concepts, virginity etc
All culty things to me.
Interviewer: Sorry also just to confirm, are "people dressed in white"  a cult or some kind of centric attribute to these "cult" groups or phenomenon you mentioned after?
Interviewee: What comes to my mind when I hear the word cult is people dressed in white
Interviewee: I took it as the first image that comes to my mind
Interviewee: My most inmediate angle of knowledge is visual, so it's what I visualized
Interviewee: Happy to answer further

---
Interview Notes: gathered on instagram after posting a story asking for people to share their thoughts on the word "cult." The interview was conducted via direct message, on Thursday, August 13th 2026. Personal information have been anonymized.


\subsection{Interview ig-05}\label{int:ig-05}
\subsubsection*{Metadata}
\begin{itemize}
  \item \textbf{ID}: ig-05
  \item \textbf{Date/time}: 2026-08-13 (date only — no time of day given in source)
  \item \textbf{Method}: Instagram
  \item \textbf{Language}: English
  \item \textbf{Location}: Germany
  \item \textbf{Interviewee}: age 34, female, Ukrainian nationality, main language Ukrainian
  \item \textbf{Questions / Answers}: 2 / 2
  \item \textbf{Total word count}: 111
  \item \textbf{Corrections log}: \texttt{cleaned/ig-05/corrections.log}
  \item \textbf{Notes}: Instagram DM interview; shortest in the corpus (single exchange).
\end{itemize}
\subsubsection*{Transcript}
Age: 34,Gender: female, Nationality: Ukrainian, Main language: Ukrainian, Language of interview: English, lives in Germany
---
Interviewer: When you hear the word cult, what comes to your mind?
Interviewee: New age
Interviewer: Amazing thank you. Briefly would you mind telling me why you associate them with a cult?
Interviewee: Because they are an enclosed group with one strange thing in the full center of their attention (for each of the new age groups)

---
Interview Notes: gathered on instagram after posting a story asking for people to share their thoughts on the word "cult." The interview was conducted via direct message, on Thursday, August 13th 2026. Personal information have been anonymized.


\subsection{Interview ig-06}\label{int:ig-06}
\subsubsection*{Metadata}
\begin{itemize}
  \item \textbf{ID}: ig-06
  \item \textbf{Date/time}: 2026-08-13 (date only — no time of day given in source)
  \item \textbf{Method}: Instagram
  \item \textbf{Language}: English
  \item \textbf{Location}: unknown — redacted in source
  \item \textbf{Interviewee}: age redacted, redacted, redacted nationality, main language unknown — redacted in source
  \item \textbf{Questions / Answers}: 2 / 2
  \item \textbf{Total word count}: 122
  \item \textbf{Corrections log}: \texttt{cleaned/ig-06/corrections.log}
  \item \textbf{Notes}: Demographics explicitly redacted by the source per its own ethical note (see thesis/04\_Appendix/3\_Appendix.tex, Ethical Notice).
\end{itemize}
\subsubsection*{Transcript}
 Language of interview: English
 Notes: Demographics redacted. See ethical note
---
Interviewer: When you hear the word cult, what comes to your mind?
Interviewee: Design \& Computation [study programme MA Design and Computation at the University of the Arts Berlin.]
Interviewer: Amazing thank you. Briefly would you mind telling me why you associate them with a cult?
Interviewee: attempting to deviate from academic norms through “breaking boundaries”, central figure with authoritative behavior and high controlled dynamic with attempts at isolating members from larger communities

---
Interview Notes: gathered on instagram after posting a story asking for people to share their thoughts on the word "cult." The interview was conducted via direct message, on Thursday, August 13th 2026. Personal information have been anonymized.


\subsection{Interview ig-07}\label{int:ig-07}
\subsubsection*{Metadata}
\begin{itemize}
  \item \textbf{ID}: ig-07
  \item \textbf{Date/time}: 2026-08-13 (date only — no time of day given in source)
  \item \textbf{Method}: Instagram
  \item \textbf{Language}: English
  \item \textbf{Location}: Germany
  \item \textbf{Interviewee}: age 32, male, Dutch / Italian nationality, main language Dutch
  \item \textbf{Questions / Answers}: 4 / 4
  \item \textbf{Total word count}: 273
  \item \textbf{Corrections log}: \texttt{cleaned/ig-07/corrections.log}
\end{itemize}
\subsubsection*{Transcript}
Age: 32, Gender: Male, Nationality: Dutch / Italian, Main language: Dutch, Language of interview: English, lives in Germany
---
Interviewer: When you hear the word cult, what comes to your mind?
Interviewee: Religion
Interviewer: Amazing thank you. Briefly would you mind telling me why you associate them with a cult?
Interviewee: Sure, there’s a certain aspect/socio-cultural norms and beliefs within groups of people that to me are central to cults. Once I started looking at religions through those norms and codes, things became clearer to me in terms of what is and is not of value to me within religions as well. Btw, I wouldn’t necessarily call all religions cults, but it is useful to compare them to create a better understanding of religions.
Interviewer: Ok great, would you have specific examples of things you would consider cults then ?
Interviewee: What do you mean with 'things I would consider cults'? Examples of these norms, codes and beliefs that feel "culty" to me?
Interviewer: Examples of things you would say " this is a cult"
Interviewee: Pfft to be honest I'm not thinking about or dealing with cults a lot. But I suppose most commonly I'd consider groups of people with some commonly considered 'alternative' beliefs and a person, figure or certain idea as a central element within their communal belief system. Also a lack of uncertainty could maybe be an interesting characteristic of cults.

---
Interview Notes: gathered on instagram after posting a story asking for people to share their thoughts on the word "cult." The interview was conducted via direct message, on Thursday, August 13th 2026. Personal information have been anonymized.


